\documentclass[11pt]{article}
\usepackage[utf8]{inputenc}
\usepackage{amsmath, amssymb}
\usepackage{booktabs}
\usepackage{geometry}
\geometry{margin=2.5cm}
\usepackage{hyperref}

\title{TPF PV-Methode A: Konvergenzanalyse der äußeren Q-Schleife\\
	\large Zusammenfassung}
\author{}
\date{}

\begin{document}
	\maketitle
	
	\section{Der Algorithmus (Methode A)}
	Die äußere Schleife aktualisiert die Blindleistung $q$ an den PV-Knoten,
	während die innere Fixpunktiteration (FPI) die Spannungen berechnet:
	\begin{equation}
		q^{(\ell+1)} = q^{(\ell)} - \omega\, D^{-1} g\big(q^{(\ell)}\big),
		\qquad
		g_k(q) = V_{\mathrm{spec},k}^2 - |V_k(q)|^2 .
	\end{equation}
	\begin{itemize}
		\item $D = \operatorname{diag}(2X_{kk})$ mit $X_{kk} = \operatorname{Im}(Z_B[k,k])$
		(Thévenin-Approximation, konstant).
		\item Innere FPI: $V_{n+1} = K\,(S^\ast \odot V^{\ast-1}) + L$ mit
		$K = -Y_{dd}^{-1}$, $L = K\,Y_{ds}\,v_s$ (beide konstant).
	\end{itemize}
	
	\section{Herleitung der Jacobi-Matrix $J_G$}
	Mit dem Update-Operator $G(q) = q - \omega D^{-1} g(q)$ gilt:
	\begin{align}
		J_G &= \frac{\partial G}{\partial q}
		= I - \omega D^{-1}\frac{\partial g}{\partial q}.
	\end{align}
	Da $V_{\mathrm{spec}}$ konstant ist, folgt
	$\dfrac{\partial g}{\partial q} = -H$ mit der wahren Sensitivitätsmatrix
	\begin{equation}
		H_{kj} = \frac{\partial |V_k|^2}{\partial q_j}.
	\end{equation}
	Einsetzen liefert:
	\begin{equation}
		\boxed{\,J_G = I + \omega\, H\, D^{-1}\,},
		\qquad \text{Konvergenz} \iff \rho(J_G) < 1.
	\end{equation}
	\textbf{Bemerkung zur Reihenfolge:} $D^{-1}H$ und $HD^{-1}$ sind ähnlich
	und besitzen dieselben Eigenwerte; für den Spektralradius $\rho(J_G)$
	ist die Reihenfolge daher unerheblich.
	
	\section{Warum Approximation $D$ statt echter Sensitivität $H$?}
	\begin{table}[h]
		\centering
		\begin{tabular}{lll}
			\toprule
			& Wahre $H$ & Thévenin $D$ \\
			\midrule
			Kosten & $n_{pv}$ volle Lastflüsse & gratis (Diagonale) \\
			Änderung pro Iteration & ja & nein (konstant) \\
			Batch-/Tensorstruktur & zerstört & erhalten \\
			Genauigkeit & exakt & gut bei $H \approx D$ \\
			\bottomrule
		\end{tabular}
	\end{table}
	\begin{itemize}
		\item $K, L, D$ werden \emph{einmalig} vorberechnet $\Rightarrow$ volle
		Parallelisierbarkeit (Kernvorteil von TPF).
		\item $H$ müsste pro äußerer Iteration neu bestimmt werden und würde die
		batched Matrixmultiplikation aufbrechen (vgl.\ Methode B).
		\item Für radiale Netze mit moderatem $R/X$ gilt $H \approx D$,
		also $\rho(J_G) < 1$.
		\item Versagt bei hohem $R/X$ (z.\,B.\ Salazar, $R/X \approx 5.82$)
		oder starker PV-PV-Kopplung.
	\end{itemize}
	
	\section{Warm-Start der inneren Schleife}
	Die Spannung $V$ wird \emph{einmal} vor der äußeren Schleife initialisiert
	und über \texttt{V\_init=V} in jede innere FPI weitergereicht:
	\begin{itemize}
		\item $\ell=0$: Flat-Start $V=1$.
		\item $\ell \geq 1$: Start von der zuvor konvergierten Lösung.
	\end{itemize}
	Nur $q$ (und damit $s_{\text{work}}$ an PV-Knoten) ändert sich zwischen
	äußeren Iterationen.
	
	\section{Beobachtungen (Netz \texttt{sz\_20\_r005}, $\omega=1.0$)}
	Parameter: $n_{bus}=19$, $n_{pv}=1$, $\eta=0.147$, $\rho=0.338$, $R/X=5.82$.
	\begin{table}[h]
		\centering
		\begin{tabular}{lcc}
			\toprule
			& Warm-Start & Cold-Start \\
			\midrule
			Äußere Iterationen        & 9 & 9 \\
			Innere Iterationen (total)& \textbf{43} & 59 \\
			Start-Residuum je Outer   & schrumpft $\approx \rho$ & konstant $\sim 1.1\cdot10^{-2}$ \\
			Innere Iter.\ je Outer    & fallend (7$\to$3) & konstant ($\sim$7) \\
			\bottomrule
		\end{tabular}
	\end{table}
	Kernaussagen:
	\begin{enumerate}
		\item Die Zahl der \textbf{äußeren} Iterationen ist unabhängig vom
		Start und durch $\rho(J_G)=0.338$ bestimmt.
		\item Warm-Start spart $\sim 27\%$ innere Iterationen.
		\item Die Start-Residuen der inneren FPI schrumpfen beim Warm-Start
		geometrisch mit dem Faktor $\approx \rho$:
		\begin{equation}
			\frac{1.18\text{e-}3}{9.36\text{e-}3}\approx 0.13,\quad
			\frac{3.50\text{e-}4}{1.18\text{e-}3}\approx 0.30,\quad
			\frac{1.14\text{e-}4}{3.50\text{e-}4}\approx 0.33 \to \rho.
		\end{equation}
		Das ist ein empirischer Fingerabdruck des Spektralradius.
		\item Beim Cold-Start ($V=1$ pro Outer) geht dieses Signal verloren.
	\end{enumerate}
	
	\section{Fazit}
	\begin{itemize}
		\item Die äußere Konvergenz wird durch $\rho(J_G) = \rho(I + \omega H D^{-1})$
		bestimmt.
		\item Die Thévenin-Approximation $D$ erhält die Tensor-Performance,
		setzt aber $H \approx D$ voraus.
		\item Warm-Start beschleunigt nur die innere Arbeit, nicht die Anzahl
		äußerer Q-Korrekturen.
	\end{itemize}
	
\end{document}